\documentclass[11pt]{article}
\usepackage[a4paper,margin=26mm]{geometry}
\usepackage{amsmath,amssymb,amsthm}
\usepackage[T1]{fontenc}
\usepackage{booktabs}
\usepackage{array}
\usepackage{upquote}
\usepackage{xurl}
\usepackage{hyperref}
\hypersetup{colorlinks=true,linkcolor=blue,urlcolor=blue,citecolor=blue,
 pdftitle={The weight spectrum of the Reed-Muller code RM(7,14) up to four weights},
 pdfauthor={Horace Dong}}
\newtheorem{lemma}{Lemma}[section]
\newtheorem{proposition}[lemma]{Proposition}
\newtheorem*{thmA}{Theorem A}
\newtheorem*{thmB}{Theorem B}
\newtheorem*{thmC}{Theorem C}
\newtheorem*{propCp}{Proposition C$'$}
\newtheorem*{propD}{Proposition D}
\newtheorem*{propE}{Proposition E}
\newtheorem*{propF}{Proposition F}
\theoremstyle{remark}
\newtheorem{remark}[lemma]{Remark}
\newcommand{\F}{\mathbb F}
\newcommand{\RM}{\mathrm{RM}}
\DeclareMathOperator{\wt}{wt}
\DeclareMathOperator{\supp}{supp}
\DeclareMathOperator{\spn}{span}
\newcommand{\cat}{\mathbin{\|}}
\newcommand{\sa}{\mathsf a}
\newcommand{\sbb}{\mathsf b}
\setlength{\emergencystretch}{1.5em}
\interfootnotelinepenalty=10000
\DeclareUrlCommand\code{\urlstyle{tt}\def\UrlBreaks{\do\_\do\/\do\.}\def\UrlBigBreaks{}}
\title{The weight spectrum of the Reed--Muller code $\mathrm{RM}(7,14)$\\ up to four weights}
\author{Horace Dong\\ \small Independent researcher\\
\small\texttt{maxwelldhx@gmail.com}\\
\small ORCID: \href{https://orcid.org/0009-0002-7554-7548}{0009-0002-7554-7548}}
\date{September 2026}
\begin{document}
\maketitle
\begin{abstract}
The binary Reed--Muller code $\RM(7,14)$ has minimum distance $d=128$.
Leuenberger and Albrizzio (arXiv:2606.21425) showed that every even number from
$312$ to $2^{14}-312$ is a weight of $\RM(7,14)$, except possibly $22$ values.
We show that $14$ of them are weights, by explicit codewords. We also give
explicit codewords, checked by two independent programs, for all even numbers
from $320$ to $2^{14}-320$ except $322$, $326$, $330$, $334$ and their
complements. Hence an open question of Carlet on the weight spectra of
$\RM(m-c,m)$ has, for $(c,m)=(7,14)$, a positive answer if and only if these
four numbers are weights. Together with the theorem of Kasami, Tokura and Azumi
on weights below $2.5d$, which we cite but could not consult, this determines
the weight spectrum of $\RM(7,14)$ except for these four numbers and their
complements; the equivalence with Carlet's question does not depend on that
theorem. Conjecture~2 of Lou and Wang, which would imply a positive answer, is
false for $m=6$. For the open weights we prove non-existence in restricted
families and reduce each of them to a question about $\RM(6,13)$ and
$\RM(7,13)$; calibrated searches found no codeword of weight $322$, $326$,
$330$ or $334$, which is heuristic evidence only. The $17$ explicit weights
are formally verified in Lean~4. The results were found by AI coding agents working under the author's
direction; their role is described at the end of the paper.
\end{abstract}

\section{Introduction}

For $0\le r\le m$, the binary Reed--Muller code $\RM(r,m)$ consists of the
truth tables on $\F_2^m$ of the polynomials in $x_1,\dots,x_m$ over $\F_2$ of
degree at most $r$, where $x_i^2=x_i$. We identify a codeword with its
polynomial (its algebraic normal form), write $\wt f$ for the number of points
at which $f=1$, and call the set of weights of the codewords the \emph{weight
spectrum}. The minimum distance of $\RM(r,m)$ is $d=2^{m-r}$. Kasami and
Tokura~\cite{kt} determined the weights below $2d$, and Kasami, Tokura and
Azumi~\cite{kta} those below $2.5d$. The whole weight spectrum of $\RM(m-c,m)$
is known, for all $m\ge2c$, for $c\le2$ (classical), for $c=3,4$ (Carlet and
Sol\'e~\cite{cs}), for $c=5$ (Carlet~\cite{c24}) and for $c=6$ (Lou and
Wang~\cite{lw}).
Carlet asked~\cite{c24} whether, for every $c\ge1$ and $m\ge2c$, the weight
spectrum of $\RM(m-c,m)$ consists of $0$, $2^m$, the weights in $[2^c,2^{c+1}]$
given by Kasami and Tokura and in $[2^{c+1},2^{c+1}+2^{c-1}]$ given by Kasami,
Tokura and Azumi, their complements to $2^m$, and \emph{all} even numbers in
$[2^{c+1}+2^{c-1},2^m-2^{c+1}-2^{c-1}]$. (We paraphrase the wording of Carlet's
extended abstract~\cite{cdmd} of~\cite{c24}, which is also that
of~\cite[Conjecture~1]{la}; we could not access the full text of~\cite{c24}.
In the earlier conjecture of Carlet and Sol\'e~\cite[Section~6]{cs} the last
interval starts at $3\cdot2^c$.) Since $d=2^c$, the last interval starts at
$2.5d$. The first open case is $(c,m)=(7,14)$: here $d=128$, and the question
asks whether every even number from $320$ to $2^{14}-320=16064$ is a weight of
$\RM(7,14)$.

Leuenberger and Albrizzio~\cite[Proposition~4]{la} showed that every even
number from $312$ to $2^{14}-312$ is a weight of $\RM(7,14)$, except possibly
the $22$ elements of
\begin{multline*}
M=\{322,326,330,334,354,8174,8178,8182,8186,8188,8190,\\
8194,8196,8198,8202,8206,8210,16030,16050,16054,16058,16062\};
\end{multline*}
parts of their proof rest on computer searches whose details are not given
(Lemmas~6 and~9 of~\cite{la}). They were unsure whether this is the whole
spectrum, and remarked that their constructions would produce all of $M$ from
the eight weights $322,326,330,334,354,4378,4380,4382$. Their constructions
follow Lou and Wang, whose Conjecture~2~\cite{lw} would give $322$, $326$,
$330$, $334$ and $354$ (see Theorem~C).

\paragraph{Results.}
Throughout, we write
\begin{gather*}
U_0=\{322,326,330,334\},\qquad U=U_0\cup(2^{14}-U_0)=U_0\cup\{16050,16054,16058,16062\},\\
K=\{0,128,192,224,240,248,252,254,256,272,288,296,304,308,312,314,316,318\};
\end{gather*}
by the theorem of Kasami, Tokura and Azumi, $K$ is the set of weights of
$\RM(7,14)$ below $320$; see (P4a) in Section~\ref{sec:prelim}.

\begin{thmA}
The following $17$ numbers are weights of $\RM(7,14)$: $354$, $4378$, $4380$,
$4382$, $8174$, $8178$, $8182$, $8186$, $8188$, $8190$, $8194$, $8196$,
$8198$, $8202$, $8206$, $8210$ and $16030$. Table~\ref{tab:wit} gives
codewords with at most seven monomials.
\end{thmA}

Fourteen of these numbers, the elements of $M\setminus U$, are listed as
missing in~\cite{la}; to our knowledge they are new. The other three, $4378$,
$4380$ and $4382$, are not new: they are weights by~\cite[Proposition~4]{la},
and the remark of~\cite{la} lists them among the eight weights from which its
constructions would produce $M$. We include codewords for them because we do
not use~\cite[Proposition~4]{la}.

\begin{thmB}
Let $S_0=\bigl(K\cup\{320,322,\dots,16064\}\cup(2^{14}-K)\bigr)\setminus U$.
The weight spectrum of $\RM(7,14)$ is $S_0\cup X$ for a set $X\subseteq U$ with
$X=2^{14}-X$. Every element of $S_0$ is the weight of an explicit codeword with
at most eight monomials.
\end{thmB}

In Carlet's question the weights below $2.5d$ and their complements are, by
definition, those given by~\cite{kt,kta}, and all weights are even by (P1). So
for $(c,m)=(7,14)$ the question asks exactly whether every even number from
$320$ to $16064$ is a weight. By the explicit codewords of Theorem~B and by
(P1), \emph{it has a positive answer if and only if $322$, $326$, $330$ and
$334$ are weights}. This equivalence does not depend on the theorem of Kasami,
Tokura and Azumi; only the determination of the whole spectrum in Theorem~B
does, where that theorem is the only external input: it excludes weights below
$320$ outside $K$ (see (P4a)). We could not consult~\cite{kta} ourselves.

Construction~2 of~\cite{lw} takes
$g_1=x_1\cdots x_m+x_{i_1}\cdots x_{i_m}+x_{i_{m+1}}\cdots x_{i_{2m}}$ and
$g_2=x_1\cdots x_m+x_{i_{2m+1}}\cdots x_{i_{3m}}+x_{i_{3m+1}}\cdots x_{i_{4m}}$
with $1\le i_j\le2m$, and $a_1,a_2\in\F_2$, and forms the concatenation
$g=0\cat(g_1+a_1)\cat g_2\cat(g_1+g_2+a_2)\in\RM(m+1,2m+2)$ of four truth
tables (Section~\ref{sec:lw}). Conjecture~2 of~\cite{lw} states that, for
$m\ge4$, the weights of these codewords include $2^{m+2}+2^m+2i$ and
$2^{2m}+2^m+2i$ for $0\le i<2^m$; for $m=6$ the first family is
$\{320,322,\dots,446\}$.

\begin{thmC}
Conjecture~2 of~\cite{lw} is false for $m=6$: no codeword of Construction~2
with $m=6$ has weight $322$. This remains true if $g_1$ and $g_2$ range over all
of $\RM(6,12)$, as in the proof of~\cite[Proposition~5]{lw}.
\end{thmC}

These codewords even miss all of $U$ (Proposition~C$'$), and so do the sums of
at most three indicators of flats of dimension at least $7$ (Proposition~D) and
the sums of at most five monomials of degree at most $7$ (Proposition~E, by
exhaustive computation). Proposition~F reduces each $w\in U_0$ to a question
about $\RM(6,13)$ and $\RM(7,13)$. Searches based on it found no codeword of
weight in $U_0$, although on $\RM(6,12)$ they find the analogues
$162,166,170,174$ of $U_0$ within two minutes; this is heuristic evidence only
(Section~\ref{sec:search}). Theorem~A is formally verified in Lean~4
(Section~\ref{sec:lean}).

\paragraph{Prior work.}
We found no earlier claim that an element of $M\setminus U$ is a weight of
$\RM(7,14)$, no determination of its spectrum, and no disproof of Conjecture~2
of~\cite{lw}. We searched on 26 September 2026, 06:41--07:15 and 13:43--14:03
UTC: the arXiv (listing of~\cite{la}, still at version~1 at 13:43 UTC, and API
searches); the works citing~\cite{la,lw,c24,cs} in Semantic Scholar; Crossref;
zbMATH Open; the MathSciNet reference lookup; MathDB, whose entries on
Carlet's question and on
Conjecture~2 of~\cite{lw} were open, without solutions or comments;
MathOverflow and Mathematics Stack Exchange; the IACR ePrint archive; HAL;
GitHub, including \texttt{trureturing}~\cite{tru} and other repositories of
AI-assisted work on open problems; and web search. Citation queries to OpenAlex
and keyword searches of Semantic Scholar hit rate limits; the full texts
of~\cite{kta}, \cite{c24} and the published version of~\cite{lw} were not
accessible, so we quote~\cite{lw} and~\cite{la} from their arXiv versions~1;
we did not search Google Scholar, and Scopus was not accessible. One possibly
relevant paper could not be read: the note~\cite{lw2} of Lou and Wang (online
4 April 2026), whose title refers to two conjectures about the weight spectra
of Reed--Muller codes; no abstract was available, its reference list in
Crossref includes~\cite{c24,cs,kt,kta,lw}, and~\cite{la} does not cite it. It
may concern the conjectures of~\cite{lw}, and so may anticipate
Theorem~C or other results here. A final check on 26 September 2026 between
15:49 and 15:52 UTC (the arXiv, where~\cite{la} still had only version~1;
GitHub; MathDB; and MathOverflow) found nothing new. This reports the coverage
of our searches, not a guarantee of priority.

\section{Preliminaries}\label{sec:prelim}
A \emph{$k$-flat} is an affine subspace of dimension $k$, with $2^k$ points; its
indicator $1_F$ is a product of $m-k$ affine functions, so it lies in
$\RM(m-k,m)$. A nonempty intersection of flats is a flat. We write
$x_S=\prod_{i\in S}x_i$ ($x_\emptyset=1$). Affine bijections preserve degrees,
weights and flats. A \emph{quadratic function} has degree at most $2$. We use
the following known results.

{\setlength{\leftmargini}{3.2em}%
\begin{itemize}
\item[(P1)] If $r<m$, every codeword of $\RM(r,m)$ has even weight, since
$\wt f\bmod2$ is the coefficient of $x_1\cdots x_m$ in $f$. As
$\wt(f+1)=2^m-\wt f$, the weight spectrum is invariant under $w\mapsto2^m-w$.
\item[(P2)] The weights of $\RM(r,m)$ are divisible by
$2^{\lfloor(m-1)/r\rfloor}$ (McEliece~\cite{mc}). In particular the weights of
$\RM(6,13)$ are divisible by $4$.
\item[(P3)] The codewords of weight $d$ are the indicators of $(m-r)$-flats
\cite[Chapter~13]{ms}. By the theorem of Kasami and Tokura~\cite{kt}, in the
form stated in~\cite[Section~2]{caims}, every codeword of $\RM(r,m)$, $r\ge2$,
with weight in $[d,2d)$ is affinely equivalent to
$x_1\cdots x_{r-2}(x_{r-1}x_r+x_{r+1}x_{r+2}+\dots+x_{r+2l-3}x_{r+2l-2})$ with
$2\le2l\le m-r+2$, or to
$x_1\cdots x_{r-l}(x_{r-l+1}\cdots x_r+x_{r+1}\cdots x_{r+l})$ with
$3\le l\le\min(r,m-r)$; both have weight $2d-2^{m-r+1-l}$. Hence in
$\RM(6,12)$ ($d=64$) a codeword of weight $96$ is $1_Gq$ with $G$ an $8$-flat
and $q$ quadratic; one of weight $112$ is of this form or is $1_F+1_{F'}$ with
$6$-flats $F,F'$, $|F\cap F'|=8$; and one of weight $126$ is $1_F+1_{F'}$ with
$6$-flats $F,F'$, $|F\cap F'|=1$.
\item[(P4a)] The weights of $\RM(7,14)$ below $320$ are exactly the elements
of $K$, by the theorem of Kasami, Tokura and
Azumi~\cite[pp.~392\,ff.]{kta}.\footnote{We could not consult~\cite{kta}
ourselves; the page reference is the one given by
Carlet~\cite[Section~1]{cdmd} and by Carlet and Sol\'e~\cite[Section~6]{cs}.
No source we read gives the complete list of the weights of $\RM(7,14)$
between $256$ and $320$; Carlet~\cite[Section~2]{caims} says that the weights
between $2d$ and $2.5d$ are too numerous to be recalled. Leuenberger and
Albrizzio~\cite[Proposition~4]{la} assert only that the elements of $K$ are
weights, not that there are no others below $320$. The part of (P4a) below
$2d=256$ follows from (P3), and the computation after Proposition~E agrees
with (P4a).}
\item[(P4b)] The weights of $\RM(6,12)$ below $160$ are $0$, $64$, $96$,
$112$, $120$, $124$, $126$, $128$, $136$, $144$, $148$, $152$, $154$, $156$
and $158$ (\cite[Lemma~2]{lw}, which cites~\cite{kt,kta}).
\end{itemize}}

\begin{lemma}\label{lem:ie}
For $S_1,\dots,S_k\subseteq\{1,\dots,m\}$,
\[
\wt\Bigl(\sum_{i=1}^kx_{S_i}\Bigr)=\sum_{\emptyset\ne T\subseteq\{1,\dots,k\}}
(-2)^{|T|-1}\,2^{m-|\bigcup_{i\in T}S_i|}.
\]
In particular, the weight depends only on the sizes of the $2^k-1$ Venn regions
of $S_1,\dots,S_k$.
\end{lemma}
\begin{proof}
For $y_1,\dots,y_k\in\{0,1\}$ the exclusive or of the $y_i$ equals
$\sum_{\emptyset\ne T}(-2)^{|T|-1}\prod_{i\in T}y_i$, by induction on $k$ from
$y\oplus y'=y+y'-2yy'$. Put $y_i=x_{S_i}(x)$ and sum over $x\in\F_2^m$: the
product over $T$ equals $1$ exactly when $x_j=1$ for all
$j\in\bigcup_{i\in T}S_i$, which happens at $2^{m-|\bigcup_{i\in T}S_i|}$
points. The size of a union is a sum of sizes of Venn regions.
\end{proof}

\begin{lemma}\label{lem:halves}
Let $f\colon\F_2^{m+1}\to\F_2$, and put $g(x)=f(x,0)$ and $g'(x)=f(x,1)$ for
$x\in\F_2^m$. Then $f=g+x_{m+1}(g+g')$, and $f\in\RM(r+1,m+1)$ if and only if
$g\in\RM(r+1,m)$ and $g+g'\in\RM(r,m)$. Moreover
$\wt f=\wt g+\wt g'=\wt(g+g')+2|\supp g\cap\supp g'|$, and
$\supp g\cap\supp g'=\supp g\setminus\supp(g+g')$.
\end{lemma}
\begin{proof}
The identity holds at $x_{m+1}=0$ and at $x_{m+1}=1$. The polynomials $g$ and
$x_{m+1}(g+g')$ have no monomial in common, so $\deg f\le r+1$ if and only if
$\deg g\le r+1$ and $\deg(g+g')\le r$. The rest is counting.
\end{proof}

\section{Explicit codewords: Theorems A and B}\label{sec:AB}

\begin{table}[t]\centering\footnotesize
\begin{tabular}{@{}r>{\raggedright\arraybackslash}p{0.66\textwidth}l@{}}
\toprule
$w$ & codeword $f_w$ & $(\sigma_1,\dots,\sigma_k)$\\
\midrule
$354$ & $\sa+\sbb+x_1x_2x_3x_5x_6x_7x_{11}+x_1x_3x_6x_7x_{10}x_{12}x_{13}+x_3x_4x_6x_7x_{10}x_{12}x_{13}$ & $(640,203,50,12,1)$\\
$4378$ & $\sa+\sbb+x_3x_4x_6+x_1x_5x_7x_{11}x_{13}+x_6x_{14}$ & $(6912,1537,173,21,1)$\\
$4380$ & $x_2x_6x_7x_9x_{12}+x_3x_4x_5x_8x_9x_{10}x_{13}+x_8x_{14}+x_1x_3x_5x_9x_{10}x_{11}x_{14}$ & $(4864,288,25,1)$\\
$4382$ & $\sa+\sbb+x_1x_2x_5x_6x_7+x_4x_9+x_3x_4x_6x_8x_{12}$ & $(5376,725,150,20,1)$\\
$8174$ & $\sa+\sbb+x_7+x_5x_{10}+x_2x_3x_4x_5x_8x_{10}x_{11}+x_1x_2x_4x_6x_{13}x_{14}$ & $(12928,2825,314,59,9,1)$\\
$8178$ & $\sa+\sbb+x_7+x_2x_5x_6x_{13}+x_2x_5x_8x_{10}x_{13}x_{14}$ & $(9728,1081,187,19,1)$\\
$8182$ & $\sa+\sbb+x_8+x_1x_2x_3x_4x_6x_{11}x_{13}+x_1x_{14}+x_2x_7x_8x_{10}x_{11}x_{12}x_{14}$ & $(12800,2765,304,50,8,1)$\\
$8186$ & $\sa+\sbb+x_{11}+x_4x_{12}+x_1x_2x_8x_{10}x_{11}x_{12}x_{13}$ & $(12672,2597,219,23,1)$\\
$8188$ & $x_1+x_2x_5x_6x_8x_9x_{10}x_{13}+x_2x_{14}+x_1x_4x_6x_7x_{11}x_{12}x_{14}$ & $(12544,2370,100,2)$\\
$8190$ & $\sa+\sbb+x_2+x_1x_7x_8$ & $(10496,1313,82,1)$\\
\bottomrule
\end{tabular}
\caption{The codewords of Theorem~A; $\sa=x_1\cdots x_7$, $\sbb=x_8\cdots
x_{14}$. If $f_w=\sum_{i=1}^kx_{S_i}$, then $\sigma_j=\sum_{|T|=j}
2^{14-|\bigcup_{i\in T}S_i|}$ over $T\subseteq\{1,\dots,k\}$, and $\wt
f_w=\sum_j(-2)^{j-1}\sigma_j$ (Lemma~\ref{lem:ie}). The complements $1+f_w$
for $w=354,8190,8188,8186,8182,8178,8174$ have the weights
$16030,8194,8196,8198,8202,8206,8210$.}\label{tab:wit}
\end{table}

\begin{proof}[Proof of Theorem~A]
Each $f_w$ in Table~\ref{tab:wit} is a sum of distinct monomials of degree at
most $7$, so $f_w\in\RM(7,14)$. By Lemma~\ref{lem:ie}, $\wt
f_w=\sum_j(-2)^{j-1}\sigma_j$ with the $\sigma_j$ of the table; for instance
$\wt f_{354}=640-2\cdot203+4\cdot50-8\cdot12+16=354$. The other seven weights
are those of the complements $1+f_w$, by (P1). All $17$ weights were also
checked by machine (Sections~\ref{sec:comp} and~\ref{sec:lean}).
\end{proof}

\begin{remark}\label{rem:direct}
The weights $8174$ to $8210$ also follow from a direct sum: if $A$ depends only
on $x_1,\dots,x_7$ and $B$ only on $x_8,\dots,x_{14}$, then $\deg(A+B)\le7$ and
$\wt(A+B)=\wt A(128-\wt B)+(128-\wt A)\wt B=8192-2(64-\wt A)(64-\wt B)$. As
$\wt A$ and $\wt B$ can take any value from $0$ to $128$, taking $\wt B=63$ and
$\wt A=64-n$ shows that every even number $8192-2n$ with $|n|\le64$ is a
weight.
\end{remark}

\begin{proof}[Proof of Theorem~B]
The file \code{rm714_all_witnesses.json} of~\cite{repo} lists, for each of the
$7901$ elements $w$ of $S_0$, a sum of at most eight distinct monomials of
degree at most $7$ with weight $w$, as two independent programs confirm
(Section~\ref{sec:comp}). So every element of $S_0$ is a weight. Conversely,
by (P1) every weight is even; by (P4a) the weights below $320$ lie in $K$, and
by complementation the weights above $16064$ lie in $2^{14}-K$. So the spectrum lies in $S_0\cup U$, and it equals $S_0\cup X$ with
$X$ the set of weights in $U$; $X=2^{14}-X$ by (P1).
\end{proof}

Most entries of the file were found by random sampling of sparse polynomials,
and $362$ by~\cite[Lemma~8]{la}; this plays no role in the proof. As $S_0$
consists of the weights asserted in~\cite[Proposition~4]{la} and the new
weights of Theorem~A, the file also checks that proposition independently.



\section{Construction 2 of Lou and Wang: Theorem C}\label{sec:lw}
Following~\cite{lw}, for functions $f_1,f_2,f_3,f_4$ on $\F_2^{12}$ let
$f_1\cat f_2\cat f_3\cat f_4$ be the function on $\F_2^{14}$ equal to $f_1$,
$f_2$, $f_3$, $f_4$ on the cosets $(x_{13},x_{14})=(0,0),(1,0),(0,1),(1,1)$, so
that its weight is $\sum\wt f_i$. For $A,B\in\RM(6,12)$ we have
$0\cat A\cat B\cat(A+B)=x_{13}A+x_{14}B\in\RM(7,14)$.

\begin{lemma}\label{lem:triples}
Let $A,B\in\RM(6,12)$ and $f=x_{13}A+x_{14}B$. Then
$\wt f=\wt A+\wt B+\wt(A+B)$. If $\wt f\in U_0$, then $\wt f\ne322$, and the
multiset $\{\wt A,\wt B,\wt(A+B)\}$ is $\{64,126,136\}$ if $\wt f=326$,
$\{64,112,154\}$ if $\wt f=330$, and one of $\{64,112,158\}$,
$\{64,126,144\}$, $\{96,112,126\}$ if $\wt f=334$.
\end{lemma}
\begin{proof}
The first claim is the weight of the concatenation. Each of $A$, $B$, $A+B$ is
the sum of the other two, and $\supp(Y+Z)\subseteq\supp Y\cup\supp Z$, so each
of the three weights is at most the sum of the other two. Hence, writing
$a\le b\le c$ for the three weights, $c\le a+b$ and $c\le\wt f/2\le167$. By (P1)
and (P4b), $a,b,c$ lie in
\[
L=\{0,64,96,112,120,124,126,128,136,144,148,152,154,156,158,160,162,164,166\}.
\]
The triples in $L$ with $a\le b\le c\le a+b$ and $a+b+c\in U_0$ are the ones
listed (certificate \code{cert_arith}, Section~\ref{sec:comp}). For
$\wt f=322$ this is also easy to see by hand. Then $c\le161$, so $a,b,c\le160$.
The elements of $L$ that are $2$ modulo $4$ are $126$, $154$ and $158$, and
$322\equiv2\pmod4$, so an odd number of $a,b,c$ are among them; as three of them
sum to at least $378$, exactly one is. But
none of $322-126=196$, $322-154=168$ and $322-158=164$ is a sum of two elements
of $\{0,64,96,112,120,124,128,136,144,148,152,156,160\}$.
\end{proof}

\begin{proof}[Proof of Theorem~C]
In both readings $g_1,g_2\in\RM(6,12)$, so let $g_1,g_2\in\RM(6,12)$ be
arbitrary and $g=0\cat(g_1+a_1)\cat g_2\cat(g_1+g_2+a_2)$. If $a_1=a_2=0$, then
$g=x_{13}g_1+x_{14}g_2$, and $\wt g\ne322$ by Lemma~\ref{lem:triples}. If
$a_1=a_2=1$, then $g=0\cat G_1\cat g_2\cat(G_1+g_2)=x_{13}G_1+x_{14}g_2$ with
$G_1=g_1+1\in\RM(6,12)$, and the same holds. If $a_1=1$ and $a_2=0$, then
$\wt g=(4096-\wt g_1)+\wt g_2+\wt(g_1+g_2)\ge4096$, as
$\wt g_2+\wt(g_1+g_2)\ge\wt g_1$; similarly $\wt g\ge4096$ if $a_1=0$ and
$a_2=1$. So $322$ is not a weight of Construction~2 for $m=6$, whereas
Conjecture~2 of~\cite{lw} asserts that $2^8+2^6+2=322$ is one.
\end{proof}

\begin{propCp}
Let $A,B\in\RM(6,12)$. Then $\wt(x_{13}A+x_{14}B)\notin U$. Consequently no
codeword of Construction~2 with $m=6$ has weight in $U$, whether $g_1$ and
$g_2$ are as in~\cite{lw} or arbitrary in $\RM(6,12)$.
\end{propCp}
\begin{proof}
We use the fact $(\ast)$: if $G,E$ are flats of $\F_2^{12}$ and $q$ is
quadratic, then $|G\cap E\cap\supp q|$ is even or at most $3$, since $q$
restricted to the flat $G\cap E$ (if nonempty, of dimension $k$) has even weight
if $k\ge3$ by (P1), and weight at most $2^k\le4$ otherwise.

The weight of $f=x_{13}A+x_{14}B$ is at most $3\cdot2^{12}=12288<16050$, and it
is not $322$ by Lemma~\ref{lem:triples}. For the five remaining triples of
Lemma~\ref{lem:triples} we choose two of the functions $A$, $B$, $A+B$, say
$X$ and $Y$, so that $X+Y$ is the third; then
$|\supp X\cap\supp Y|=(\wt X+\wt Y-\wt(X+Y))/2$. We use (P3) throughout.

\emph{The triples $\{64,126,136\}$ and $\{64,126,144\}$.} Take $X=1_E$ with $E$
a $6$-flat and $Y=1_F+1_{F'}$ with $F\cap F'=\{p\}$. Then
$|E\cap\supp Y|=27$, resp.\ $23$. As $\supp Y=(F\cup F')\setminus\{p\}$,
$|E\cap F|+|E\cap F'|=|E\cap\supp Y|+2[p\in E]$ is $27$ or $29$, resp.\ $23$ or
$25$. Both summands are $0$ or powers of $2$, so one of them is $1$, and the
other is $22$, $24$, $26$ or $28$, which is impossible.

\emph{The triples $\{64,112,154\}$ and $\{64,112,158\}$.} Take $X=1_E$ with
$E$ a $6$-flat and $\wt Y=112$. Then $|E\cap\supp Y|=11$, resp.\ $9$. If
$Y=1_Gq$, this contradicts $(\ast)$. Otherwise $Y=1_F+1_{F'}$ with
$|F\cap F'|=8$. Put $s=|E\cap F|$, $s'=|E\cap F'|$ and $t=|E\cap F\cap F'|$,
each $0$ or a power of $2$. Then $s+s'-2t$ is $11$, resp.\ $9$. It is odd, so
$s=1$ or $s'=1$; by symmetry let $s=1$. Then $t\le1$, and $s'=10+2t$,
resp.\ $s'=8+2t$, so we are in the second case with $s'=8$ and $t=0$. By (P3)
we may take $F=\{x_1=\dots=x_6=1\}$ and $F'=\{x_1=x_2=x_3=x_7=x_8=x_9=1\}$
(the form $x_1x_2x_3(x_4x_5x_6+x_7x_8x_9)$). Since $|E\cap F|=1$, the affine
map $E\to\F_2^6$, $x\mapsto(x_1,\dots,x_6)$, has a fibre with one point, so it
is injective and hence bijective. So the set $E_3$ of points of $E$ with
$x_1=x_2=x_3=1$ has $8$ points. As $E\cap F'\subseteq E_3$ and
$|E\cap F'|=8$, we get $E\cap F'=E_3$. The point of $E\cap F$ lies in $E_3$,
hence in $F'$, so $t=1$, a contradiction.

\emph{The triple $\{96,112,126\}$.} Take $X=1_Gq$ of weight $96$ and
$Y=1_F+1_{F'}$ of weight $126$ with $F\cap F'=\{p\}$, so that $\wt(X+Y)=112$.
Then $|\supp X\cap\supp Y|=55$, and $|\supp X\cap F|+|\supp X\cap F'|$ is $55$
or $57$. One summand is odd, hence at most $3$ by $(\ast)$. The other is at most
$48$. Indeed, it is the weight of $q$ on the flat $G\cap F$ (or $G\cap F'$). If
this flat has dimension at most $5$, the weight is at most $32$. Otherwise
$F\subseteq G$, and $q$ is not constant $1$ on $F$: otherwise
$\supp X\supseteq F$, and $X+1_F$ would be a nonzero codeword of $\RM(6,12)$
of weight $32<64$. So $1+q$ restricted to $F$ is a nonzero codeword of
$\RM(2,6)$, of weight at least $16$, and $q$ has weight at most $48$ on $F$.
The total is at most $51<55$, a contradiction.

Finally, in Construction~2 the cases $a_1=a_2$ give codewords
$x_{13}A+x_{14}B$ with $A,B\in\RM(6,12)$, as in the proof of Theorem~C, and the
other cases give weights between $4096$ and $12288$.
\end{proof}

In particular, the codewords of~\cite[Constructions~1--3, Propositions~1
and~2]{la}, which are of this form, have no weight in $U$.

\section{Partial non-existence results for the open weights}\label{sec:nonex}

\subsection{Sums of flats}
Call two flats \emph{transversal} if they meet in exactly one point. Two
transversal flats of $\F_2^{14}$ of dimension at least $7$ both have dimension
$7$: if their direction spaces are $V$ and $V'$, then $V\cap V'=\{0\}$, so
$\dim V+\dim V'\le14$.

\begin{propD}
Let $F_1,F_2,F_3$ be flats of $\F_2^{14}$ of dimension at least $7$, not
necessarily distinct (the whole space is allowed). Then none of $\wt 1_{F_1}$,
$\wt(1_{F_1}+1_{F_2})$ and $\wt(1_{F_1}+1_{F_2}+1_{F_3})$ lies in $U$.
\end{propD}
\begin{proof}
The weight of $1_{F_1}+1_{F_2}+1_{F_3}$ is
$\sum|F_i|-2\sum_{i<j}|F_i\cap F_j|+4|F_1\cap F_2\cap F_3|$, and a nonempty
intersection of flats has a power of $2$ as its size. As $4\mid|F_i|$, the
weight is $2$ modulo $4$ if and only if an odd number of the three pairs are
transversal. The same holds for one or two flats. All elements of $U$ are $2$
modulo $4$.

For one flat there is no transversal pair. Two transversal $7$-flats give the
weight $128+128-2=254$. Three pairwise transversal $7$-flats give
$384-6+4|F_1\cap F_2\cap F_3|\in\{378,382\}$. It remains to treat exactly one
transversal pair, say $F_1\cap F_2=\{p\}$ with $\dim F_1=\dim F_2=7$, and
$\dim F_3=d\in[7,14]$. Translate so that $p=0$; then $F_1=V_1$ and $F_2=V_2$
are linear subspaces with $V_1\oplus V_2=\F_2^{14}$. Let $V_3$ be the direction
space of $F_3$. The pairs $(F_1,F_3)$ and $(F_2,F_3)$ are not transversal.

\emph{Case $p\in F_3$.} Then $F_3=V_3$, $|V_1\cap V_3|=2^a$ and
$|V_2\cap V_3|=2^b$ with $a,b\ge1$, and the weight is
$258+2^d-2^{a+1}-2^{b+1}$. Since $(V_3\cap V_1)\oplus(V_3\cap V_2)\subseteq V_3$,
we have $a+b\le d$. The only solutions of $258+2^d-2^{a+1}-2^{b+1}\in U$ with
$1\le a\le b\le7$ and $7\le d\le14$ are $(w,d,a,b)=(322,7,4,4)$ and
$(322,8,5,6)$, and both violate $a+b\le d$.

\emph{Case $p\notin F_3$.} Then the weight is $254+2^d-2(\alpha+\beta)$, where
$\alpha=|F_1\cap F_3|$ and $\beta=|F_2\cap F_3|$ are $0$ or powers of $2$ from
$2$ to $128$. The
only solution in $U$ is $w=334$, $d=7$, $\{\alpha,\beta\}=\{8,16\}$. Then
$\dim(V_3\cap V_1)+\dim(V_3\cap V_2)=3+4=7=\dim V_3$, so
$V_3=(V_3\cap V_1)\oplus(V_3\cap V_2)$. Write $F_3=c+V_3$ and $c=c_1+c_2$ with
$c_i\in V_i$. As $F_3$ meets $F_1$, $c\in V_1+V_3=V_1+(V_3\cap V_2)$, so
$c_2\in V_3\cap V_2$; similarly $c_1\in V_3\cap V_1$. Hence $c\in V_3$ and
$0=p\in F_3$, a contradiction.

Both statements about solutions are checked by \code{cert_arith}, and by hand
they are quick, as $2^{a+1}+2^{b+1}$ and $\alpha+\beta$ have at most two binary
digits $1$.
\end{proof}

The obstruction is dimension counting: in $\RM(6,12)$ the $6$-flats spanned by
$e_1,\dots,e_6$, by $e_7,\dots,e_{12}$ and by $e_1,e_2,e_3,e_7,e_8,e_9$ give
$162=2.5\cdot64+2$, whereas $322$ in $\RM(7,14)$ would need $a=b=4$, against
$a+b\le7$.

\subsection{Sums of few monomials}
\begin{propE}
No sum of at most five monomials of degree at most $7$ in $x_1,\dots,x_{14}$
(the constant $1$ allowed) has weight in $U$. The same holds for $f\circ\phi$ for
every such sum $f$ and every affine bijection $\phi$ of $\F_2^{14}$.
\end{propE}
\begin{proof}[Proof by exhaustive computation]
By Lemma~\ref{lem:ie}, the weight of $\sum_{i=1}^kx_{S_i}$ depends only on the
Venn region sizes $n_T$, $\emptyset\ne T\subseteq\{1,\dots,k\}$, which satisfy
$\sum_Tn_T\le14$ and $|S_i|=\sum_{T\ni i}n_T\le7$. Equal monomials cancel in
pairs, so the sums of at most five distinct monomials are exactly the sums
$\sum_{i=1}^kx_{S_i}$ with $k\in\{4,5\}$ and not necessarily distinct $S_i$. We
enumerated all vectors $(n_T)$ for $k\le5$: there are $8$, $204$, $20{,}708$,
$8{,}568{,}777$ and $11{,}690{,}447{,}828$ of them for $k=1,\dots,5$. Two
independent programs (Section~\ref{sec:comp}) agree on these counts, and
neither finds a weight in $U$. The second statement holds because $\phi$
preserves weights.
\end{proof}

Among these sums, the weights $2$ modulo $4$ in $[250,400]$ are $254$, $314$,
$318$ and all such numbers from $338$ to $398$, and the weights below $320$ are
exactly the elements of $K$, in agreement with (P4a).

\subsection{A reduction}
\begin{lemma}\label{lem:dir}
Let $f\in\RM(7,14)$ have support $S$ and weight $w\equiv2\pmod4$. For
$a\in\F_2^{14}\setminus\{0\}$ let $n_a$ be the number of pairs $\{x,x+a\}$
contained in $S$. Then $\sum_{a\ne0}n_a=\binom w2$, every $n_a$ is odd, and in
coordinates in which $a$ is the last unit vector we have $f=g+x_{14}h$ with
$g\in\RM(7,13)$, $h\in\RM(6,13)$, $n_a=|\supp g\setminus\supp h|$ and
$\wt h=w-2n_a$.
\end{lemma}
\begin{proof}
Each $2$-subset $\{x,y\}$ of $S$ is counted once, by $a=x+y$. Composing $f$
with an invertible linear map sending $e_{14}$ to $a$, we may assume
$a=e_{14}$. By Lemma~\ref{lem:halves} with
$m=13$ and $r=6$, $f=g+x_{14}h$ with $g=f(\cdot,0)\in\RM(7,13)$ and
$h=f(\cdot,0)+f(\cdot,1)\in\RM(6,13)$; the pairs in direction $e_{14}$
correspond to the points of $\supp g\cap\supp(g+h)=\supp g\setminus\supp h$,
and $\wt h=w-2n_a$. By (P2), $4\mid\wt h$, so $2n_a\equiv w\equiv2\pmod4$.
\end{proof}

\begin{propF}
Let $w\in U_0$. Then $w$ is a weight of $\RM(7,14)$ if and only if there are
$h\in\RM(6,13)$ and $g\in\RM(7,13)$ with
\[
(\wt h,\ |\supp g\setminus\supp h|)=(w-2,1)\quad\text{or}\quad(w-6,3).
\]
Moreover, a codeword of weight $w=322,326,330,334$ has such a decomposition, as
in Lemma~\ref{lem:dir}, in at least $7559$, $7235$, $6908$, $6576$ directions.
\end{propF}
\begin{proof}
If $g$ and $h$ are as stated, then $f=g+x_{14}h\in\RM(7,14)$ has weight
$\wt h+2|\supp g\setminus\supp h|=w$ by Lemma~\ref{lem:halves}. Conversely,
let $f$ have weight $w$. By Lemma~\ref{lem:dir} the $n_a$ are odd and
$\sum_{a\ne0}(n_a-1)=\binom w2-16383$. A direction with $n_a\ge5$ contributes
at least $4$ to this sum, so at most $\lfloor(\binom w2-16383)/4\rfloor$
directions have $n_a\ge5$; for $w=322,326,330,334$ this bound is $8824$,
$9148$, $9475$, $9807$. In each of the remaining directions $n_a\in\{1,3\}$, and
Lemma~\ref{lem:dir} gives the decomposition with $\wt h=w-2n_a$.
\end{proof}

\section{Computations}\label{sec:comp}

\subsection{Certificates}
Every computation used in a proof was done by two programs written separately
(in \code{certificates/rm714/} of version~1.4.0 of~\cite{repo}, with data and
logs). For Theorems~A and~B, \code{verify_witnesses.c} (binary M\"obius
transform) and \code{verify_bitset.py} with \code{verify_all_py.py}
(big-integer bit sets) check degrees and weights; both accept the codewords of
Table~\ref{tab:wit}, their complements and the $7901$ entries of
\code{rm714_all_witnesses.json} (SHA-256 \texttt{784bca3b}\dots), whose
weights are exactly the elements of $S_0$. For Lemma~\ref{lem:triples} and
Propositions~C$'$, D and~F, \code{cert_arith.py} and \code{cert_arith.c} print
the same $38$ lines. For Proposition~E, \code{venn_enum.c} (for $k\le4$) and
\code{venn_enum_fast.c} (for $k=5$) recurse over the Venn regions and use
Lemma~\ref{lem:ie}, while \code{venn_enum_indep.c} splits the regions monomial
by monomial and uses M\"obius inversion. Their logs agree on the counts for
every $k$, on all weights up to $400$ for $k\le4$, and, for $k=5$, on the
weights $2$ modulo $4$ in $[250,400]$ and on the number $5676$ of distinct
weights. No run finds a weight in $U$. The complete sets of weights, which
\code{venn_enum_indep.c} writes to files, were not compared.

\subsection{Searches for the open weights (heuristic)}\label{sec:search}
For fixed $h\in\RM(6,13)$, the existence of $g\in\RM(7,13)$ with
$|\supp g\setminus\supp h|=n$, for $n=1$ or $3$ as in Proposition~F, is linear
algebra: let $H=\supp h$ and let
$c(x)\in\F_2^{2380}$ be the values at $x\in\F_2^{13}$ of the monomials of
degree at most $5$, a column of a parity-check matrix of $\RM(7,13)$, whose dual
is $\RM(5,13)$~\cite[Chapter~13]{ms}. A vector supported on a set $P$ disjoint
from $H$ agrees outside $H$ with a codeword if and only if
$\sum_{x\in P}c(x)\in\spn c(H)$. The program \code{closure_search.c} samples
structured $h$ of degree at most $6$ (sums of two to four pieces, each the
indicator of a $7$- to $9$-flat or the indicator of a $9$- to $11$-flat times
a sum of one to three products of two or three affine forms; or, with the
option \texttt{GEN\_FLATS}, sums of three or four such pieces, mostly
indicators of $7$- and $8$-flats) and decides $n=1$ and $n=3$ exactly for each
$h$ of a relevant weight. It first maps the $c(x)$ by a fixed random linear map
to $\F_2^{512}$, which can create false positives but no false negatives; each
positive is re-solved with the full vectors and the codeword is rebuilt and
rechecked. No positive was false, and \code{check_closure_hits.py} rechecks in
pure Python every codeword printed in the logs ($18$ distinct ones for each
code). With $m=11$ the same program runs on $\RM(6,12)$ as a calibration.

\begin{table}[t]\centering\footnotesize\setlength{\tabcolsep}{4pt}
\begin{tabular}{@{}rrrrr@{\hspace{14pt}}rrrrr@{}}
\toprule
\multicolumn{5}{c}{$\RM(7,14)$, $h\in\RM(6,13)$} &
\multicolumn{5}{c}{$\RM(6,12)$, $h\in\RM(5,11)$}\\
\cmidrule(r){1-5}\cmidrule(l){6-10}
& \multicolumn{2}{c}{$n=1$} & \multicolumn{2}{c}{$n=3$}
& & \multicolumn{2}{c}{$n=1$} & \multicolumn{2}{c}{$n=3$}\\
$w$ & sampled & found & sampled & found & $w$ & sampled & found & sampled & found\\
\midrule
$318^*$ & $2782+549$ & $2782+549$ & $4285+1010$ & $1360+276$ & $158^*$ & $3420$ & $3245$ & $4581$ & $1616$\\
$322$ & $6574+1958$ & $0$ & $2782+549$ & $0$ & $162$ & $9581$ & $58$ & $3420$ & $0$\\
$326$ & $15+89$ & $0$ & $6574+1958$ & $0$ & $166$ & $1213$ & $366$ & $9581$ & $58$\\
$330$ & $229+728$ & $0$ & $15+89$ & $0$ & $170$ & $3996$ & $1678$ & $1213$ & $803$\\
$334$ & $170+587$ & $0$ & $229+728$ & $0$ & $174$ & $3052$ & $2820$ & $3996$ & $1679$\\
$338^*$ & $1159+3623$ & $43+165$ & $170+587$ & $0$ & $178^*$ & $32029$ & $886$ & $3052$ & $754$\\
$342^*$ & $1059+2140$ & $239+700$ & $1159+3623$ & $43+165$ & & & & &\\
\bottomrule
\end{tabular}
\caption{Searches of Section~\ref{sec:search}: for a target weight $w$ and
$n=1$ (resp.\ $3$), the number of distinct sampled $h$ of weight $w-2$
(resp.\ $w-6$), and the number of them admitting $g$ with
$|\supp g\setminus\supp h|=n$, i.e.\ giving a codeword of weight $w$. Entries
$x+y$: default generator $+$ \texttt{GEN\_FLATS} (one $h$ of weight $340$ was
sampled by both). Starred: known weights, used as controls.}\label{tab:search}
\end{table}

The runs for $\RM(7,14)$ (26 September 2026, 13:59--14:33 UTC) were five of
$1200$ seconds with the default generator (seeds $11$--$15$), five of $720$
seconds with \texttt{GEN\_FLATS} (seeds $41$--$45$) and one of $240$ seconds
(seed $51$); the calibration was one run of $120$ seconds (seed $21$).
Consecutive seeds give shifted copies of one random stream, and the runs of each
group sampled the same sequence of $h$. Regenerating the samples from the
recorded seeds and iteration counts (\code{closure_replay.c}) confirms this, and
shows that no $h$ of these weights repeats within a run and that the run with
seed $51$ sampled a subset of the first group. Table~\ref{tab:search} counts each $h$ once:
$26{,}956$ distinct $h$ of weights $312$ to $340$, of which $13{,}681$ have the
weights $316$ to $332$ relevant for $U_0$.

In $\RM(6,12)$ all four analogues $162,166,170,174$ of $U_0$ were found, and
$162=2.5\cdot64+2$ arose only from $h$ of weight exactly $2.5\cdot64$, with
$n=1$ ($58$ of $9581$ such $h$). In $\RM(7,14)$ none of the $8532$ sampled $h$
of weight $320=2.5\cdot128$ admits $n=1$, and none of the $3331$ of weight
$316$ admits $n=3$. Earlier, cruder searches found no weight in $U_0$ either
(they did not reach weights near $2^{14}-U_0$): random sparse polynomials
(\code{sparsefind.c}; the weights it found are listed in
\code{sparse_table_s7.txt}), and about $6.7\cdot10^9$ random sums of three to
six flat indicators and hill-climbing steps (\code{flatsearch.c}), which met
$318$ and $338$ millions of times. All this is heuristic: the test is exact for
each sampled $h$, but the sampled $h$ are few and special.

\section{Formal verification}\label{sec:lean}
Theorem~A is formalised in \code{RM714Weights.lean} and
\code{RM714WeightsAudit.lean} (in \code{lean/Research/} of version~1.4.0
of~\cite{repo}; Lean~4.33.1, Mathlib~\cite{mathlib} commit \texttt{0df444a3}).
In the namespace \texttt{RM714} a polynomial is a list of monomials, a monomial
being the set $S\subseteq\{0,\dots,13\}$ of the indices of its variables
(index $i$ stands for $x_{i+1}$), and \texttt{evalPoly L x} is
\texttt{(L.map fun S => prod i in S, x i).sum} in \texttt{ZMod 2}. In ASCII
notation (\texttt{\#} is the cardinality of a finite set) the other definitions
and the main theorem read:
\begin{footnotesize}
\begin{verbatim}
def weight (L : List (Finset (Fin 14))) : Nat :=
  #{x : Fin 14 -> ZMod 2 | evalPoly L x = 1}
def IsRM714Weight (w : Nat) : Prop :=
  exists L : List (Finset (Fin 14)), (forall S in L, S.card <= 7) /\ weight L = w
theorem rm714_new_weights :
  forall w in [354, 16030, 4378, 4380, 4382, 8174, 8178, 8182, 8186, 8188,
               8190, 8194, 8196, 8198, 8202, 8206, 8210], IsRM714Weight w
\end{verbatim}
\end{footnotesize}
So \texttt{IsRM714Weight w} says exactly that $w$ is a weight of $\RM(7,14)$
(repeated monomials are allowed; they cancel).

\emph{A note on the name.} Despite its name, \texttt{rm714\_new\_weights}
certifies all $17$ weights of Theorem~A, of which only $14$ are new: $354$,
$8174$, $8178$, $8182$, $8186$, $8188$, $8190$, $8194$, $8196$, $8198$, $8202$,
$8206$, $8210$ and $16030$, the elements of $M\setminus U$. The other three,
$4378$, $4380$ and $4382$, are asserted by~\cite[Proposition~4]{la}. The name
is historical; the claim of novelty rests on the text of this paper, not on the
name of the declaration.

The weights are computed on $14$-bit masks: \texttt{weightNat} counts the bit
patterns at which the exclusive or of the monomials is $1$;
\texttt{weight\_map\_maskSet} proves that this count is \texttt{weight} (using
Mathlib's \texttt{finFunctionFinEquiv}); and \texttt{isRM714Weight\_of} turns
masks of degree (\texttt{degNat}) at most $7$ with count $w$ into
\texttt{IsRM714Weight w}. The witnesses \texttt{L354}, \dots, \texttt{L8210}
(Table~\ref{tab:wit} and complements) are checked by
\texttt{isRM714Weight\_354}, \dots, \texttt{isRM714Weight\_8210} with
\texttt{decide +kernel}, that is, by evaluation in the kernel. An audit file
prints the definitions and statements and tests the definitions on examples.
There is no \texttt{sorry}, no \texttt{axiom} and no \texttt{native\_decide};
\texttt{\#print axioms} for \texttt{rm714\_new\_weights},
\texttt{weight\_map\_maskSet} and \texttt{isRM714Weight\_of} reports only
\texttt{propext}, \texttt{Classical.choice} and \texttt{Quot.sound}. A rebuild
with \texttt{lake build} in a clean directory (fresh sources, Mathlib build from
the local cache) and \texttt{leanchecker} on both modules exited with code $0$
(26 September 2026, 13:49--13:55 UTC; the log is in version~1.4.0 of the
repository). Not formalised are Theorem~B (its $7901$ codewords would need
hours of kernel time in this encoding, and its upper bound rests
on~\cite{kta}), Theorem~C, Propositions~C$'$, D, E and~F, and the searches.

\section{Open questions}\label{sec:open}
Are $322$, $326$, $330$ and $334$ weights of $\RM(7,14)$? This is all that
remains of its spectrum, and it decides Carlet's question for $(c,m)=(7,14)$.
For $322$, Proposition~F suggests a route: the case $n=3$ concerns the
codewords of weight $316$ of $\RM(6,13)$, described in~\cite{kta}, and for each
normal form the existence of $g$ is finite linear algebra; the case $n=1$ asks
whether some $h\in\RM(6,13)$ of weight exactly $320$ has a point
$p\notin\supp h$ with $c(p)\in\spn c(\supp h)$, as some codewords of weight
$160$ of $\RM(5,11)$ do in the calibration.

\section*{Tool and computational resource disclosure}
This work was carried out on 26 September 2026 (UTC) in a workflow directed by
the author, using AI coding agents: Anthropic Claude Code, with the model
Claude Opus~5.5. Claude Code agents selected the problem, found the results and
their proofs, wrote the programs and the Lean formalisation, carried out the
literature searches described in the introduction, and drafted this text; the
programs, data, logs and Lean code are in~\cite{repo}. Throughout, ``we''
refers to this workflow; the reading and searching described in the paper were
done by the agents. The computations ran on one computer (macOS, arm64) on the
same day: Proposition~E took about $21$ CPU minutes with
\code{venn_enum_fast.c} and about $8$ minutes on four processes with
\code{venn_enum_indep.c}, the runs behind Table~\ref{tab:search} about $2.8$
CPU hours, those of \code{flatsearch.c} $2$ hours, and the clean Lean rebuild
with \texttt{leanchecker} about $5$ minutes. The author is not a professional
mathematician.
The author has read the paper and, using a side-by-side table, compared the
statements of the Lean theorems with the results stated here.
The results rest on the written proofs, on the codewords of
Section~\ref{sec:AB} with their machine checks, on the computation behind
Proposition~E and on the known results (P1)--(P4b); the searches of
Section~\ref{sec:search} are heuristic and not used in any proof. AI systems are not authors of this paper; the author
takes full responsibility for its content.

{\footnotesize
\begin{thebibliography}{99}
\bibitem{cdmd}
C.~Carlet, The weight spectrum of the Reed--Muller codes $\RM(m-5,m)$,
extended abstract, Discrete Mathematics Days 2024,
\url{https://dmd2024.web.uah.es/files/abstracts/paper_10.pdf}.
\bibitem{c24}
C.~Carlet, The weight spectrum of the Reed--Muller codes $\RM(m-5,m)$,
\emph{IEEE Trans.\ Inf.\ Theory} \textbf{70}(7) (2024), 4799--4807,
\url{https://doi.org/10.1109/TIT.2023.3343697}.
\bibitem{caims}
C.~Carlet, Identifying codewords in general Reed--Muller codes and determining
their weights, \emph{AIMS Math.} \textbf{9}(5) (2024), 10609--10637,
\url{https://doi.org/10.3934/math.2024518}.
\bibitem{cs}
C.~Carlet and P.~Sol\'e, The weight spectrum of two families of Reed--Muller
codes, \emph{Discrete Math.} \textbf{346}(10) (2023), 113568,
\url{https://doi.org/10.1016/j.disc.2023.113568}.
\bibitem{repo}
H.~Dong, \emph{ai4math-results}, repository,
\url{https://github.com/Horace-Maxwell/ai4math-results}, version~1.4.0
(2026); archived at Zenodo under the concept DOI
\url{https://doi.org/10.5281/zenodo.22970254}.
\bibitem{kt}
T.~Kasami and N.~Tokura, On the weight structure of Reed--Muller codes,
\emph{IEEE Trans.\ Inf.\ Theory} \textbf{16}(6) (1970), 752--759,
\url{https://doi.org/10.1109/TIT.1970.1054545}.
\bibitem{kta}
T.~Kasami, N.~Tokura and S.~Azumi, On the weight enumeration of weights less
than $2.5d$ of Reed--Muller codes, \emph{Inform.\ and Control}
\textbf{30}(4) (1976), 380--395,
\url{https://doi.org/10.1016/S0019-9958(76)90355-7}.
\bibitem{la}
M.~Leuenberger and M.~Albrizzio, On the weight spectrum of the Reed--Muller
codes $\RM(7,14)$, arXiv:2606.21425v1, 19 June 2026.
\bibitem{lw}
Y.~Lou and Q.~Wang, Determining the weight spectrum of the Reed--Muller codes
$\RM(m-6,m)$, \emph{Des.\ Codes Cryptogr.} \textbf{93} (2025), 4925--4936,
\url{https://doi.org/10.1007/s10623-025-01708-7}; quoted from
arXiv:2406.03803v1.
\bibitem{lw2}
Y.~Lou and Q.~Wang, A note on two conjectures about the weight spectra of the
Reed--Muller codes, \emph{Discrete Appl.\ Math.} \textbf{388} (2026),
142--145, \url{https://doi.org/10.1016/j.dam.2026.03.038}.
\bibitem{ms}
F.~J. MacWilliams and N.~J.~A. Sloane, \emph{The Theory of Error-Correcting
Codes}, North-Holland, Amsterdam, 1977.
\bibitem{mathlib}
The mathlib Community, The Lean mathematical library, in \emph{CPP 2020},
ACM, 2020, \url{https://doi.org/10.1145/3372885.3373824}.
\bibitem{mc}
R.~J. McEliece, Weight congruences for $p$-ary cyclic codes, \emph{Discrete
Math.} \textbf{3} (1972), 177--192.
\bibitem{tru}
The Omega Institute, \texttt{trureturing} repository,
\url{https://github.com/the-omega-institute/trureturing}, accessed 26
September 2026.
\end{thebibliography}
}
\end{document}
